\section{Robot platforms in \texttt{emet}}
\label{app:robots}

Dynagraph experiments use the same ZMQ observation protocol and DynaMem/GraphEQA stack across embodiments; camera layout, calibration, kinematics, actuator coupling and supported skills remain embodiment-specific.
Table~\ref{tab:robots} summarizes the robots planned for MuJoCo and (where noted) hardware evaluation.

\begin{table}[h]
  \centering
  \caption{Emet robot embodiments for Dynagraph evaluation.}
  \label{tab:robots}
  \small
  \begin{tabular}{@{}p{.14\linewidth}p{.20\linewidth}p{.32\linewidth}p{.25\linewidth}@{}}
    \toprule
    ID & Platform & Sim serve & Depth in sim \\
    \midrule
    \texttt{stretch} & Hello Robot Stretch & \texttt{emet serve mujoco} & Rendered sensor depth \\
    \texttt{rby1} / \texttt{galaxea\_r1} & Galaxea R1 (MolmoSpaces) & \texttt{emet serve molmospaces} & Rendered sensor depth \\
    \texttt{innate\_mars} & Innate Maurice (Mars) & \texttt{emet serve mujoco --robot innate\_mars} & Rendered stereo; see App.~\ref{app:innate_mars} \\
    \texttt{sourccey} & Sourccey MkV & \texttt{emet serve mujoco --robot sourccey} & Rendered depth; native TAMP unqualified \\
    \bottomrule
  \end{tabular}
\end{table}

\textbf{Stretch} is the default table-scene robot and shares its hardware family with GraphEQA's real-robot demonstrations.
MuJoCo merges \texttt{stretch.xml} with shared scene assets; planar base navigation uses the standard \texttt{xyt} ZMQ interface.

\textbf{Galaxea R1} (\texttt{--robot rby1}) extends evaluation to a humanoid-style platform in MolmoSpaces iTHOR and Robocasa merges.
Spawn metadata and merged MJCF paths live beside \texttt{galaxea\_r1.xml}; see project docs under \texttt{docs/molmospaces.md}.

\textbf{Innate Mars} (Maurice) is documented in Appendix~\ref{app:innate_mars}.
It is the primary case where \textbf{Depth Anything 3} (Appendix~\ref{app:da3}) replaces missing hardware depth while sim runs should keep sensor depth for reproducibility.

All robots publish the same ZMQ keys where applicable: \texttt{rgb}, optional \texttt{depth}, \texttt{camera\_K}, \texttt{camera\_pose}, \texttt{gps}, \texttt{compass}, plus stereo/tertiary fields when the MJCF or ROS bridge exposes them.
Dynagraph and GraphEQA clients consume this stream identically; navigation configuration includes dynav YAML (\texttt{depth\_source}, DA3 presets) and camera count. Shared key names do not imply equivalent manipulation controllers or sensor availability. Native TAMP acceptance in this paper is limited to stationary Stretch; the corrected Sourccey simulator model is an implementation asset, not evidence of Sourccey native task success.
